\section{Ámbito de Aplicación o Influencia}
 
\subsection{Descripción del ámbito}
 
El proyecto está orientado al entorno administrativo y académico de la \textbf{Universidad Nacional de San Antonio Abad del Cusco (UNSAAC)}, específicamente en los procesos relacionados con la gestión de los procedimientos establecidos en el Texto Único de Procedimientos Administrativos (TUPA).
 
La propuesta se desarrolla dentro del contexto de la \textbf{transformación digital universitaria}, buscando modernizar los servicios administrativos mediante el uso de tecnologías web que permitan automatizar procesos, mejorar la atención a los usuarios y optimizar el control y seguimiento de expedientes administrativos.
 
\subsection{Ámbito institucional}
 
El proyecto tiene como ámbito de aplicación directo la UNSAAC, con sede principal en Cusco y filiales en:
 
\begin{itemize}[noitemsep]
    \item Cusco (Sede Central)
    \item Sicuani
    \item Espinar
    \item Puerto Maldonado
    \item Santo Tomás
    \item Yanaoca
    \item Andahuaylas
\end{itemize}
 
El sistema PLADDES (\textit{tramite.unsaac.edu.pe}) es la plataforma oficial de trámite documentario virtual de la institución y es el objeto de mejora de este proyecto.
 
\subsection{Área de influencia}
 
Entre las principales unidades involucradas se encuentran:
 
\begin{itemize}[noitemsep]
    \item Rectorado
    \item Vicerrectorado Académico
    \item Decanatos de Facultad
    \item Escuela de Posgrado
    \item Centro de Cómputo
    \item Oficina de Admisión
    \item Mesa de Partes Virtual
    \item Unidades de Trámite Documentario
\end{itemize}
 
\subsection{Usuarios objetivo / Beneficiarios}
 
\begin{itemize}
    \item \textbf{Estudiantes de pregrado y posgrado:} principal grupo que registra solicitudes en PLADDES para trámites como certificados, constancias, homologaciones, grado y título, entre otros.
    \item \textbf{Docentes y personal cesante/administrativo:} usuarios que gestionan trámites propios de su condición laboral a través del sistema.
    \item \textbf{Dependencias administrativas:} Unidad de Trámite Documentario, Decanatos, Escuelas Profesionales, VRAC, DIGA, UTH, SG y demás oficinas que reciben y resuelven los expedientes ingresados en PLADDES.
    \item \textbf{Público en general e instituciones:} personas externas a la comunidad universitaria que acceden a servicios puntuales del TUPA.
    \item \textbf{La institución universitaria en general,} que se verá favorecida con el fortalecimiento de la modernización administrativa y la mejora en la calidad del servicio.
\end{itemize}
 
\subsection{Alcance funcional}
 
El proyecto rediseña y optimiza los flujos existentes en PLADDES, abarcando:
 
\begin{enumerate}
    \item \textbf{Registro de solicitudes:} el wizard de 4 pasos actual será reemplazado por un flujo moderno, guiado e intuitivo.
    \item \textbf{Catálogo de trámites TUPA:} presentación clara de los 80+ procedimientos con requisitos, costos, plazos y canales de atención, integrada al flujo de registro.
    \item \textbf{Seguimiento de expedientes:} mejora del módulo de consulta de estado de trámite con vista de línea de tiempo y notificaciones.
    \item \textbf{Gestión de pagos:} integración mejorada del proceso de verificación de voucher de pago, actualmente manual.
    \item \textbf{Panel administrativo:} rediseño del panel de las dependencias para gestión eficiente de la bandeja de expedientes.
    \item \textbf{Reportes y estadísticas:} módulo nuevo para generación de reportes de trámites procesados por dependencia y período.
\end{enumerate}